% Página
\section{Estudo}

A segurança de sistemas computacionais está intimamente ligada à robustez das senhas utilizadas. Este estudo foi desenvolvido com o objetivo de avaliar diferentes características de senhas, dividindo a análise em duas partes principais: senhas com propriedades fixas e senhas com propriedades variadas.

\subsection{Fixa}

Nesta etapa, foram avaliadas senhas que seguem um padrão predefinido de propriedades, incluindo:\\

\textbf{Calcular o número de combinações para cada tipo de caractere.}\\
Calculando a quantidade de caracteres possíveis.
\begin{itemize}
	\item 03 Letras maiúsculas : $26^3$ caracteres
	\item 03 Letras minúsculas : $26^3$ caracteres
	\item 03 Números : $10^3$ caracteres
	\item 03 Caracteres especiais : $8^3$ caracteres
	\item Ordem dos caracteres não importa;
	\item Os mesmos caracteres podem ser usados mais de uma vez;
\end{itemize}

\textbf{Determinar o número de arranjos possíveis}\\
Após determinar os grupos de caracteres, precisamos calcular de quantas maneiras podemos organizar os 12 caracteres (3 maiúsculas, 3 minúsculas, 3 números, e 3 caracteres especiais). Esse é um problema de permutações de multiconjuntos, onde os elementos são indistinguíveis dentro de cada grupo.\\

A fórmula geral para permutações de multiconjuntos é:

\[
P = \frac{n!}{k_1! \cdot k_2! \cdot k_3! \cdot k_4!}
\]

Onde:
\begin{itemize}
    \item $n$ é o número total de elementos,
    \item $k_1, k_2, k_3, k_4$ são os tamanhos dos grupos.
\end{itemize}

Neste caso:
\begin{itemize}
    \item $n$ = 12,
    \item $k_1 = k_2 = k_3 = k_4$  = 3.
\end{itemize}

Substituindo:

\[
P = {\frac{12!}{3! \cdot 3! \cdot 3! \cdot 3!}} { = 369.600}
\]\\

Isso significa que existem 369.600 maneiras de organizar os 12 caracteres.\\

\textbf{Calcular o número total de senhas}\\

O número total de senhas possíveis é o produto das combinações para cada grupo de caracteres e o número de arranjos possíveis:\\

\[
\text{Total de Senhas} = P \cdot \big( (26^3) \cdot (26^3) \cdot (10^3) \cdot (8^3) \big)
\]

\[
\text{Total de Senhas} = 369.600 \cdot 17.576 \cdot 17.576 \cdot 1.000 \cdot 512
\]

\[
\text{Total de Senhas} = 58.457.738.654.515.200.000 = 5,84577386545152 \times 10^{19}
\]

Este é o total de senhas possíveis para as especificações dadas.

\subsection{Variada}

Nesta etapa, a senha é gerada de forma aleatória, sem considerar as características definidas para a Senha Fixa mencionadas no item anterior exceto o tamanho de 12 caracteres.\\

\textbf{Calcular o número de combinações para cada tipo de caractere.}\\
Calculando a quantidade de caracteres possíveis.
\begin{itemize}
	\item Letras maiúsculas : $26^3$ caracteres
	\item Letras minúsculas : $26^3$ caracteres
	\item Números : $10^3$ caracteres
	\item Caracteres especiais :  - $8^3$ caracteres
\end{itemize}

Portanto, o número total de caracteres disponíveis para cada posição na senha é a soma dos tamanhos desses conjuntos:\\

\[
N_{\text{total}} = 26 + 26 + 8 + 10 = 70
\]\\

\textbf{Calcular o número total de senhas}\\

Se tivermos \( n \) caracteres possíveis e uma senha de \( L \) dígitos, o número total de combinações possíveis é dado por:

\[
\text{Total de Senhas} = n^L
\]

Onde:
\begin{itemize}
    \item \( n = 70 \) (número total de caracteres disponíveis),
    \item \( L = 12 \) (comprimento da senha).
\end{itemize}

Substituímos os valores de \( n \) e \( L \) na fórmula:

\[
\text{Total de Senhas} = 70^{12}
\]

Calculando:

\[
\text{Total de Senhas} = 13.841.287.201.000.000.000.000 = 1,3841287201 \times 10^{22}
\]

Este é o total de senhas possíveis para as especificações dadas.

\pagebreak
